\section{The class-group consequence of a single $K$-group}
This section proves the exact arithmetic transport from a $K$-group to one
component in the cyclotomic tower. Put $A_r=\Cl(F_r)\{p\}$ and
$\Lambda=\Z_p[[\Gamma]]$, with $F_r$ and $\Gamma$ as above.
For every character exponent $j$ modulo $p-1$, write
$A_r^{[j]}=\epsilon_j A_r$ and $X^{[j]}=\epsilon_j X$.

\begin{proposition}\label{prop:component-control}
The norm projection induces a $\Delta$-equivariant isomorphism
\[
 X/(\gamma-1)X\ \xrightarrow{\ \sim\ } A_1.
\]
In particular $A_1^{[j]}=0$ is equivalent to $X^{[j]}=0$ for every $j$.
\end{proposition}
\begin{proof}
Let $L_r$ be the Hilbert $p$-class field of $F_r$. Base change preserves
unramified abelian $p$-extensions, so $L_rF_{r+1}\subseteq L_{r+1}$.
The fields $L_r$ are therefore nested. Put $L_\infty=\bigcup_rL_r$ and
$G=\Gal(L_\infty/F_1)$. Class field reciprocity identifies
$\Gal(L_\infty/F_\infty)$ with $X$ and restriction with the norm maps.
Choose a prime of $L_\infty$ above $p$ and its inertia group $I$ in $G$.
At every finite layer the inertia group maps isomorphically to
$\Gal(F_r/F_1)$: $L_r/F_r$ is unramified and $p$ is totally ramified
in $F_r/F_1$. Taking inverse limits gives $I\simeq\Gamma$.
Thus $G=X\rtimes I$, its closed commutator subgroup is $(\gamma-1)X$,
and its abelianization fits into
\[
 0\longrightarrow X/(\gamma-1)X\longrightarrow G^{\rm ab}
   \longrightarrow\Gamma\longrightarrow0,
\]
with the image of $I$ as a section.

The quotient $G^{\rm ab}/I$ describes an abelian pro-$p$ extension
of $F_1$ unramified at $p$; all other primes were unramified already.
Conversely $L_1$ lies in $L_\infty$. Maximality of the Hilbert
$p$-class field therefore identifies this quotient with $A_1$.
The inclusion of $X$ induces the asserted isomorphism. It is the
canonical norm projection, hence commutes with $\Delta$, independently
of the chosen inertia section.

Apply $\epsilon_j$ and topological Nakayama. For completeness, a
nonzero compact $\Lambda$-module has a nonzero finite continuous
$\Lambda$-quotient. Reducing such a quotient modulo $p$ and taking
coinvariants for the finite $p$-group through which $\Gamma$ acts
leaves a nonzero quotient: its group algebra over $\F_p$ is local
with nilpotent augmentation ideal. Thus vanishing modulo
$(p,\gamma-1)$ forces the original compact module to vanish.
This proves the stated equivalence.
\end{proof}
The inertia argument is classical; see, for example,
\cite[\S1.3, pp.~3--4]{Wuthrich}.

\begin{proposition}\label{prop:k-component}
For every odd prime $p$ and every integer $m\geq1$, there is an
isomorphism
\[
 K_{4m}(\Z)/p\ \simeq\ (A_1/pA_1)^{[-2m]}.
\]
The following three groups vanish together:
$K_{4m}(\Z)\{p\}$, $A_1^{[-2m]}$, and $X^{[-2m]}$.
\end{proposition}
\begin{proof}
The odd-primary comparison in \cite[VI, Corollary~8.3]{Weibel}
gives
\[
 K_{4m}(\Z)/p\simeq
 H^2_{\mathrm{\acute et}}(\Z[1/p],\mu_p^{\otimes(2m+1)}).
\]
Over $S'=\operatorname{Spec}\Z[\mu_p,1/p]$, Kummer theory gives
$H^2_{\mathrm{\acute et}}(S',\mu_p)\simeq\operatorname{Pic}(S')/p$.
Indeed the $p$-primary Brauer contribution vanishes: there are no
real places and a single omitted finite prime, so the sum of local
invariants is injective on that possible summand. The omitted prime
is principal, generated by $1-\zeta_p$, whence
$\operatorname{Pic}(S')=\Cl(F_1)$.

Tensor with $\mu_p^{\otimes2m}$ and descend along the finite
\'etale $\Delta$-torsor $S'\to\operatorname{Spec}\Z[1/p]$.
Since $|\Delta|=p-1$ is invertible modulo $p$, descent takes invariants:
\[
 H^2_{\mathrm{\acute et}}(\Z[1/p],\mu_p^{\otimes(2m+1)})
 \simeq\bigl((A_1/pA_1)\otimes\mu_p^{\otimes2m}\bigr)^\Delta.
\]
The twist contributes the character $\omega^{2m}$, so these invariants
select precisely the character $\omega^{-2m}$. Both $K_{4m}(\Z)$
and $A_1$ are finite. Vanishing modulo $p$ therefore means vanishing
of the $p$-primary part, and Proposition~\ref{prop:component-control}
gives the assertion about $X$.
\end{proof}
This is the modern form of Kurihara's comparison
\cite[Corollary~1.5, Proposition~3.7 and Remark~3.11]{Kurihara};
see also \cite[VI, Theorem~10.9]{Weibel}. For the trivial character,
$A_1^{[0]}=0$ because extension after norm to $\Q$ acts as $p-1$
on that component and $\Cl(\Q)=0$.

\section{The posted $K$-group input and its dependencies}
\label{sec:dependencies}

This section states the external vanishing input used in the
application and follows its arithmetic dependencies explicitly.
Mok--Zheng's Theorem~1, p.~2, states
$K_{4m}(\mathbb Z)=0$ for $1\leq m\leq30$ \cite{MZ}.
Their manuscript is posted and unrefereed. The two comparison
maps in Propositions~\ref{prop:k-component} and
\ref{prop:component-control} carry this statement to
$X^{[-2m]}=0$ at every odd prime. For $p\leq61$ these residues
include every nontrivial even character; for $p>61$ they give
thirty distinct characters. The trivial class character is
already zero by the norm map to $\mathbb Q$.

The cases $m=1,2$ are the refereed $K_4$ theorem of Rognes
and the $K_8$ theorem of Dutour Sikiri\'c--Elbaz-Vincent--Kupers--Martinet,
respectively; see \cite[\S2, pp.~2--3]{MZ} for the exact
identifications. In the remaining range, the odd primes below
$2^{31}$ are handled by Hart--Harvey--Ong's Vandiver computation
\cite[\S1, assertion~(2)]{HHO} and the $K$-theory comparison.
The two-primary computation uses Rognes--Weibel and
\cite[VI, Corollary~9.8]{Weibel}. Our Galois application uses
only odd primes.

For odd $p>2^{31}$, Mok--Zheng use a strict Selmer comparison,
an arithmetic upper bound extracted from Alp\"oge's archive,
and exact Hecke determinants. More precisely, put $e=2m$.
Their Proposition~2, pp.~2--3, compares the $p$-primary
$K_{2e}(\mathbb Z)$ obstruction with the strict Selmer group
of $(\mathbb Q_p/\mathbb Z_p)(-e)$, unramified everywhere and
locally trivial at $2$ and $p$. The upper bound used in
\S\S3--5, pp.~3--8, is the pointwise Selmer bound in the
archive, specialized to dimension $2$, the trivial character,
and coefficient ring $\mathbb Z_p$. Thus the dependency is
through that arithmetic statement in the archive. The
numerical Riemann hypothesis for zeta and Dirichlet $L$-functions
does not replace any term in this displayed Selmer comparison.

For $3\leq m\leq26$, the auxiliary fields are
$\mathbb Q(\sqrt{-3})$ and $\mathbb Q(\sqrt{-7})$.
At each Hecke probe, a nonzero Selmer class forces a nonunit
trace difference on one of the admissible classical or
unitary branches. Multiplying the branch factors at each
probe before taking their greatest common divisor is the
logical step that permits the branch to vary with the probe.
The determinant checks in \S\S6--7, pp.~8--11, eliminate
these possibilities. For $27\leq m\leq30$, the argument uses
$\mathbb Q(i)$ and an additional transfer of the archive
construction, described in \S6.5, p.~10, and \S7, p.~11.
That transfer removes one comparison input and replaces six
substantive uses in the source argument. Its public provenance
statement reports that this modified entry point has not
been fully rebuilt.

Sutherland's verification report \cite{Sutherland} records
kernel checks of the original combined archive, including
several independent kernel implementations. Those checks
concern the original archive whose hash is given there.
They do not by themselves check the additional Gaussian
transfer. No new Lean kernel execution is claimed for the
application in this paper.

\subsection{The joint-representation route}

Pintar's posted theorem, p.~1 of \cite{Pintar}, reaches
$m\leq26$ from the arithmetic construction in the same
archive. Its \S2, p.~2, uses the joint-representation
argument: a positive strict Selmer length produces one
admissible integral representation with residual
semisimplification
\[
                 1\oplus\omega^e\oplus\omega^{-e},
\]
valid simultaneously at every Galois element. The logical
map can be stated directly. An admissible local point gives
such an integral representation. If there are no such
representations, the space of local points is empty.
Evaluating the model upper bound at the identity then has
empty covering set and value zero. The Selmer-to-model
inequality forces strict Selmer length zero. This is the
contrapositive used in Pintar's determinant argument.

This route does not use Mok--Zheng's vanishing theorem as
an input. Both dependency paths do use the archive's
arithmetic construction and the $K$-theory comparison.
Pintar's stated range ends at $m=26$, so it gives no
additional route to the final four components. The version
opened here has public identifier 2611.00829v3 and internal
title-page designation ``Version 5, 23 November 2026''.

\subsection{Soul\'e's torsion bound and its exact specialization}

The comparisons above use no general-number-field
torsion-stabilizer bound from Soul\'e's 2003 Lemma~9.
The reviewed joint-representation proof uses the Selmer
and model inequalities just stated; Mok--Zheng and Pintar
do not cite that lemma in their vanishing arguments.
Here is an explicit calculation fixing the scope of the
bound, without identifying unrelated assertions.

Let $F=\mathbb Q(\sqrt5)$, put $t=(\sqrt5-1)/2$, and set
\[
                 A=\begin{pmatrix}0&-1\\1&t\end{pmatrix}
                    \in\operatorname{SL}_2(O_F).
\]
The identities $t^2+t-1=0$ and $A^2-tA+1=0$ give
$A^5=1$, while $A\ne1$. Its order is $5$, exceeding
$1+\max([F:\mathbb Q],2)=3$.
The valid dimension argument is as follows. Regard a
nonidentity order-$p$ endomorphism of $F^N$ as a rational
endomorphism of dimension $[F:\mathbb Q]N$. Its minimal
polynomial has a primitive $p$th root as a zero and hence
contains the irreducible polynomial $\Phi_p$ of degree
$p-1$. Therefore $p\leq[F:\mathbb Q]N+1$. For
$F=\mathbb Q$ this gives $p\leq N+1$, exactly the bound
used in the rational specialization. The displayed
matrix disproves the stated larger-field bound while
the minimal-polynomial argument proves this specialization.



\section{Reflection and a finite list of arithmetic exceptions}
For $1\leq m\leq30$, put $\theta_m=2+2m$ and let
\[
 \mathcal E_m=\{p\text{ prime}:p\geq2^{31},\quad
             p\mid\operatorname{num}(B_{\theta_m})\}.
\]
The nonempty sets are listed in Table~\ref{tab:reflection-exceptions}.

\begin{proposition}\label{prop:reflection-exceptions}
For every $1\leq m\leq30$ and every odd prime
$p\notin\mathcal E_m$, the component $X^{[2+2m]}$ vanishes.
In particular this vanishing holds at every odd prime for
$m=1,\ldots,15$ and for $m=21$.
\end{proposition}
\begin{proof}
We give the reflection morphism explicitly. Put $F=\Q(\mu_p)$,
$A=\Cl(F)\{p\}$, and let $B\subset F^\times$ consist of elements
whose $p$th roots generate everywhere unramified extensions of $F$.
Kummer theory and Artin reciprocity give
\[
 B/F^{\times p}\simeq\operatorname{Hom}(A,\mu_p).
\]
For $b\in B$, every valuation of $b$ is divisible by $p$.
Write $(b)=\mathfrak a^p$. The map $[b]\mapsto[\mathfrak a]$
is well-defined, is $\Delta$-equivariant, and has kernel
$(B\cap\mathcal O_F^\times)/\mathcal O_F^{\times p}$.
For an even $\theta$ with $2\leq\theta\leq p-3$, take
$\chi=\omega^{1-\theta}$. The $\chi$ part of this kernel vanishes:
the odd part of global units modulo $p$ comes from roots of unity,
which have character $\omega$, whereas $\chi\ne\omega$.
Indeed $u/c(u)$ is a root of unity for every unit $u$, since it is
an algebraic integer with all complex conjugates of absolute value one.
Thus the displayed map restricts to an injection
\[
 \operatorname{Hom}(A^{[\theta]},\mu_p)
       \lhook\joinrel\longrightarrow A^{[1-\theta]}[p].
\]
This proves the required direction of Leopoldt reflection, including
its character shift; compare \cite[Theorem~4.4.2 and
Corollary~4.4.4]{SharifiIwa}.

For completeness, the vanishing on the right follows from a precise
Stickelberger annihilator. For $b$ primitive modulo $p$, projection
of $(\sigma_b-b)\Theta$ onto the $\omega^{1-\theta}$ component gives
\[
 (\omega^{1-\theta}(b)-b)B_{1,\omega^{\theta-1}}.
\]
The first factor is a $p$-adic unit, and
$B_{1,\omega^{\theta-1}}\equiv B_\theta/\theta\pmod p$.
Consequently $p\nmid\operatorname{num}(B_\theta)$ implies
$A^{[1-\theta]}=0$ and then $A^{[\theta]}=0$.
These are the Herbrand implication and the Bernoulli congruence in
\cite[Theorem~4.5.7 and Corollary~5.4.15]{SharifiIwa}.

If $p<2^{31}$, the published Vandiver computation
\cite[\S1, assertion~(2)]{HHO} gives $A^+=0$ directly.
For $p\geq2^{31}$, every $\theta_m$ lies between $4$ and $p-3$.
If $p\notin\mathcal E_m$, the preceding injection and annihilator
give $A^{[\theta_m]}=0$. In both cases,
Proposition~\ref{prop:component-control} gives $X^{[\theta_m]}=0$.

The arithmetic table was computed from the defining rational
Bernoulli recurrence. Every numerator factorization was checked
by exact multiplication and every asserted prime was certified
by the fully factored $n-1$ Lucas criterion, recursively down to $2$.
The standard-library verifier and all certificates accompany this paper.
They show $\mathcal E_m=\varnothing$ for $m=1,\ldots,15,21$,
which proves the last assertion.
\end{proof}

\begin{table}[htbp]
\centering
\small
\begin{tabular}{rrl}
$m$ & $\theta_m$ & $\mathcal E_m$ \\
\hline
16&34&$151628697551$\\
17&36&$26315271553053477373$\\
18&38&$154210205991661$\\
19&40&$1897170067619$\\
20&42&$1520097643918070802691$\\
22&46&$67568238839737$\\
23&48&$153289748932447906241$\\
24&50&$47464429777438199$\\
25&52&$4534045619429$\\
26&54&$1120412849144121779$\\
27&56&$19088082706840550550313$\\
28&58&$6235242049,\quad37349583369104129$\\
29&60&$109317926249509865753025015237911$\\
30&62&$28765594733083851481$\\
\hline
\end{tabular}
\caption{The fifteen prime exceptions to the reflection calculation.
The sets for $m=1,\ldots,15$ and $m=21$ are empty.}
\label{tab:reflection-exceptions}
\end{table}

This result is independent of the posted $K$-group vanishing claims.
Combining the proved reflection result with the $K$-group input stated
by Mok--Zheng supplies the pair
$X^{[-2m]}=X^{[2+2m]}=0$ for $1\leq m\leq30$ and
$p\notin\mathcal E_m$. Their input is currently an unrefereed
posted result; its dependency and verification status are recorded
separately in the audit. Pintar's posted argument supplies another
route to the $K$-group input for $1\leq m\leq26$, with shared
archival dependencies. Neither source is needed for
Proposition~\ref{prop:reflection-exceptions} itself.
No vanishing of $X^{[\theta_m]}$ at the primes in
Table~\ref{tab:reflection-exceptions} is established here.

\begin{corollary}[Application of the posted $K$-theory theorem]
\label{cor:k-canonical-detection}
For an odd prime $p$, put
\[
 \Sigma_{30}(p)=\{-2m\bmod(p-1):1\leq m\leq30\}.
\]
With the vanishing theorem of \cite[Theorem~1, p.~2]{MZ} as the
external input described above, Theorem~\ref{thm:selected-all-stages}
gives at every odd prime the canonical Galois subextensions
\[
 L_{2,\Sigma_{30}(p)}\subsetneq L_{3,\Sigma_{30}(p)}
 \subsetneq\cdots\subsetneq L_{61,\Sigma_{30}(p)}
 \subseteq\Omega_p^*.
\]
For every \(2\leq n\leq61\), the group
\(\operatorname{Gal}(L_{n,\Sigma_{30}(p)}/F_\infty)=Q_{n,\Sigma_{30}(p)}\)
has nilpotency class exactly \(n\) and abelianization
\((U/D_n)\oplus T_{\Sigma_{30}(p)}\).
There are two elements whose length-\(n\) commutator is nonzero.
Its class-three quotient has second and third lower-central
graded groups
\[
 N_{\Sigma_{30}(p)}\oplus(W\widehat\otimes T_{\Sigma_{30}(p)}),
 \qquad \operatorname{Lie}^3(W)/[W,N_{\rm bad}],
\]
respectively.
\end{corollary}
\begin{proof}
Proposition~\ref{prop:k-component} transports each stated
$K_{4m}(\mathbb Z)$ vanishing to $X^{[-2m]}=0$.
These are precisely the component hypotheses of
Theorem~\ref{thm:selected-all-stages}. For \(2\leq n\leq61\),
the characters \(-2,-4,\ldots,-2\lfloor n/2\rfloor\) all belong
to \(\Sigma_{30}(p)\), with residues read modulo \(p-1\).
That theorem proves exact class and the strict transitions.
Proposition~\ref{prop:selected-class-three} gives the displayed
lower-central groups of the class-three quotient.
\end{proof}
The Bernoulli prime list in Table~\ref{tab:reflection-exceptions}
records the limits of that reflection calculation. Every prime
in the table is included in the displayed Galois-field application,
which uses the component $X^{[-2m]}$ directly.
The characteristic-zero comparison, the reflection morphism,
and the construction of $C_E$ are independent of the posted
vanishing theorem. The assertion of the corollary uses that
external theorem, with its verification status as described in
Section~\ref{sec:dependencies}. For $m\leq26$, Pintar's posted
argument gives the alternative route with the shared archive
inputs identified there. For $p<2^{31}$ the full equality
$\Omega_p^*=\Omega_p$ already follows from \cite{HHO,Note}.

\section{The first irregular prime}
This section makes the arithmetic indices and a denominator divisor
explicit at $p=37$. The even characters have exponents
\[
 0,2,4,6,8,10,12,14,16,18,20,22,24,26,28,30,32,34.
\]
The Bernoulli recurrence gives the residues
\[
\begin{array}{c|rrrrrrrrr}
 r&2&4&6&8&10&12&14&16&18\\\hline
 B_r\bmod37&31&16&15&16&4&17&32&22&31\\
\end{array}
\]
\[
\begin{array}{c|rrrrrrrr}
 r&20&22&24&26&28&30&32&34\\\hline
 B_r\bmod37&15&15&17&12&29&2&0&2.
\end{array}
\]
In particular $B_{32}=-7709321041217/510$ has $37$-adic valuation one.
The unique irregular index is $32$, and Herbrand--Ribet gives the
odd class character $1-32\equiv5\pmod{36}$.

Here is a small Vandiver certificate. Work modulo $q=149$, put $z=16$,
which has order $37$, and use the corrected cyclotomic-unit test
of \cite[\S4, pp.~4--5]{BH}, with the correction in \cite[\S5]{HHO}:
\[
 V=\prod_{c=1}^{18}
      \bigl(z^{c/2}-z^{-c/2}\bigr)^{r_c}\pmod{149},
 \qquad r_c\in\{0,\ldots,36\},\quad r_c\equiv c^4\pmod{37}.
\]
The exponents $c/2$ are read modulo $37$; the reduced outer
exponents are allowed by the same test. Exact modular
multiplication gives $V=132$ and $V^4=81\ne1$.
Since $32$ is the only irregular index, the criterion proves
$A_1^+=0$. Thus $X^{[e]}=0$ for all eighteen even characters,
and every even characteristic ideal of $X$ has generator $1$.
The four congruences for the Iwasawa invariants in \cite[\S5]{HHO}
are also small here. With
\[
 s=37^{-1}\sum_{a=1}^{18}a^{31},\qquad
 t=37^{-1}\sum_{a=1}^{18}a^{67},
\]
we obtain, modulo $37$,
\[
 2^{32}=7,\quad s=30,\quad t=31,\quad
 (2-32)s=25,\quad(1-32)t=1.
\]
They prove $\lambda_{37}=\nu_{37}=1$. In particular
$X^{[5]}$ is a free $\Z_{37}$-module of rank one, and all other
odd class components vanish.

Choose $\gamma$ with cyclotomic value $u=38$ and write $Z=\gamma-1$.
Let $\vartheta_5(Z)$ be the inverse limit of the negative
Stickelberger projections
\[
 -\frac1{37^{n+1}}
 \sum_{\substack{1\leq a<37^{n+1}\\37\nmid a}}
 a\,\omega(a)^{-5}\gamma_a^{-1},
 \qquad
 \gamma_a\in\Gamma/\Gamma^{37^n}.
\]
This specifies the power series exactly. Its principal ideal is
$\operatorname{char}(X^{[5]})$ by the main conjecture.
Projection at $n=2$ gives
\[
 \vartheta_5(Z)=a_0+a_1Z+a_2Z^2+\cdots,
 \qquad
 \begin{aligned}
 a_0&\equiv14\cdot37+33\cdot37^2&&\pmod{37^3},\\
 a_1&\equiv16+6\cdot37&&\pmod{37^2},\\
 a_2&\equiv29&&\pmod{37}.
 \end{aligned}
\]
The coefficients agree with \cite[\S1.9, p.~8]{Wuthrich} and are
reconstructed by the program cited through the workbench.
The finite-level error is divisible by $(1+Z)^{37^2}-1$, hence
vanishes modulo $37^3$ at every $Z\in37\Z_{37}$.
The displayed precision is therefore sufficient to determine the root
modulo $37^3$. Since $a_1$ is a unit, Weierstrass preparation gives
one root $r\in37\Z_{37}$ and a unit multiple of $Z-r$.
Writing $r=37a+37^2b$, the first digit gives $14+16a\equiv0$, so
$a=13$ and $14+16\cdot13=6\cdot37$. The next digit gives
\[
 6+33+6\cdot13+29\cdot13^2+16b
       \equiv23+16b\equiv0\pmod{37},
\]
so $b=24$. Put $\beta=1+r$ and $\alpha=38/\beta$. Then
\[
 r\equiv13\cdot37+24\cdot37^2,\qquad
 \alpha\equiv1+25\cdot37+20\cdot37^2\pmod{37^3}.
\]
The characteristic polynomials for $\gamma$ on the rational even
parts of $H^{\rm ab}$ are exactly
\[
 q_{32}(z)=z-\alpha,\qquad q_e(z)=1\quad(e\ne32).
\]
Here $\alpha$ is the precisely defined $37$-adic number above,
not a rational approximation substituted for a characteristic root.
The corresponding scalar divisor and its coordinate are
\[
 D(X,Y)=(1+X)(1+Y)-\alpha,\qquad
 Y=\frac{\alpha+S}{1+X}-1,
\]
so $D=S$ and $1/D$ has a simple pole with Laurent coefficient one.
For example the divisor passes through $(0,\alpha-1)$.
The dual cup character is $2-32\equiv6$ and its eigenvalue is
$38^2/\alpha=38\beta$. This scalar calculation specifies the
arithmetic denominator; the survival of a companion numerator
on that divisor is the separate cup-comparison requirement.
